Betting markets are known to underprice favourites. The favourite--longshot bias first documented at racetracks also
appears on Kalshi and Polymarket, where political markets look slightly underconfident. We ask whether this holds
equally for primaries and general elections. Using every two-candidate US contest settled on the two platforms between
November 2023 and September 2026, priced at a fixed time on election day before polls close, we fit one number, a
stretch, that measures whether favourites win more often (above 1) or less often (below 1) than their prices imply.
General elections followed the familiar pattern (5 November 2024: 1.97 on Kalshi, 1.83 on Polymarket), but Kalshi
primaries went the other way, with favourites winning far less often than priced (0.32); Polymarket's primaries were
close to calibrated. Because most general elections in the data were from 2024 and most primaries from 2026, the split
could mean either that primaries are traded differently or that the underconfidence faded over time. We pre-registered
a test that separates the two, archived it before the US midterms of 3 November 2026, and ran it on those general
elections. In this preliminary run (187 contests on Kalshi, 145 on Polymarket; 40 still unsettled), Polymarket remained
underconfident (1.66, 95\% interval 1.21 to 2.91) and consistent with 2024, the outcome that supports primaries being
traded differently. Kalshi's result (1.32, 0.88 to 2.35) was inconclusive, although on the 141 races both platforms
listed, both were underconfident to the same degree. Buying favourites priced above 70\% earned a few percent per bet after
costs, but two general-election nights cannot show whether the bias can be traded. A final version will follow once
every contest has settled.
